\documentclass[11pt,a4paper]{article}

% --- PACKAGES ---
\usepackage[utf8]{inputenc}
\usepackage[indonesian]{babel}
\usepackage{amsmath, amssymb, amsfonts}
\usepackage{geometry}
\geometry{margin=1in, headheight=14pt}
\usepackage{graphicx}
\usepackage{booktabs}
\usepackage{hyperref}
\usepackage{tcolorbox}
\tcbuselibrary{skins,breakable}
\usepackage{listings}
\usepackage{xcolor}
\usepackage{fancyhdr}

% --- STYLING & COLORS ---
\definecolor{unibblue}{RGB}{0, 51, 102}
\definecolor{unibgold}{RGB}{212, 175, 55}
\definecolor{darkgreen}{RGB}{34, 139, 34}

\hypersetup{
    colorlinks=true,
    linkcolor=unibblue,
    citecolor=darkgreen,
    urlcolor=blue
}

\pagestyle{fancy}
\fancyhf{}
\rhead{\textbf{TKE-41XX: Optimisasi untuk Teknik Elektro}}
\lhead{Catatan Kuliah Minggu ke-14}
\rfoot{Halaman \thepage}

\lstset{
  language=Python,
  basicstyle=\ttfamily\footnotesize,
  keywordstyle=\color{unibblue}\bfseries,
  commentstyle=\color{darkgreen}\itshape,
  stringstyle=\color{red!70!black},
  numbers=left,
  numberstyle=\tiny\color{gray},
  numbersep=5pt,
  breaklines=true,
  frame=single,
  rulecolor=\color{unibblue!30},
  tabsize=4
}

\title{\textbf{Catatan Kuliah Minggu ke-14:}\\Optimisasi Multi-Objektif dan Metode Hibrid dalam Teknik Elektro\\(Sintesis Bahan Ajar EE 2023--2026 dan Rekayasa Mesin)}
\author{\textbf{Ir. Novalio Daratha S.T., M.Sc., Ph.D.}\\Program Studi S1 Teknik Elektro --- Universitas Bengkulu}
\date{Semester Ganjil 2026}

\begin{document}

\maketitle

\begin{tcolorbox}[colback=unibblue!5,colframe=unibblue,title=\textbf{Ringkasan Eksekutif dan Integrasi Kurikulum}]
Dokumen ini menyajikan materi pengajaran sintesis yang menghubungkan pengalaman optimisasi Teknik Elektro tahun-tahun sebelumnya (2023--2026: Economic Dispatch, Metaheuristik GA/PSO/SCA, Desain Filter Digital FIR) dengan metode optimisasi deterministik presisi tinggi dari kurikulum Teknik Mesin (*Armijo Line Search*, *Quasi-Newton BFGS*, dan *Augmented Lagrangian Method*). Mahasiswa mempelajari pembentukan *Pareto Frontier* pada masalah *Combined Economic dan Emission Dispatch* (CEED) dan penentuan solusi kompromi terbaik (*Best Compromise Solution*).
\end{tcolorbox}

\tableofcontents
\newpage

\section{Sintesis Konseptual: Teknik Elektro dan Rekayasa Mesin}
Mata kuliah Optimisasi untuk Teknik Elektro bertujuan membekali mahasiswa S1 dengan dua pilar utama:
\begin{enumerate}
    \item **Pilar Metaheuristik (Tradisi EE)**: Fleksibilitas tinggi dalam menangani fungsi tujuan non-konveks, non-diferensiabel, dan *search space* yang luas (seperti pergerakan partikel PSO dan seleksi kromosom GA).
    \item **Pilar Deterministik Presisi (Tradisi Teknik Mesin)**: Kecepatan konvergensi lokal superlinier dan kepatuhan kendala yang sangat ketat melalui *Quasi-Newton BFGS*, *Armijo Line Search*, dan *Augmented Lagrangian Method*.
\end{enumerate}

Dengan menggabungkan kedua pilar ini, insinyur elektro dapat menyelesaikan masalah kompleks pada *Smart Grid* secara cepat, akurat, dan serempak memenuhi beberapa fungsi tujuan yang saling bertentangan.

\section{Teori Optimisasi Multi-Objektif (MOO)}

\subsection{Formulasi Umum}
Secara umum, masalah optimisasi multi-objektif dengan $M$ fungsi tujuan dan $K$ kendala dapat dituliskan sebagai:
\begin{equation}
    \min_{\mathbf{x} \in \Omega} \mathbf{F}(\mathbf{x}) = \left[ f_1(\mathbf{x}), f_2(\mathbf{x}), \dots, f_M(\mathbf{x}) \right]^T
\end{equation}
dengan kendala:
\begin{equation}
    g_j(\mathbf{x}) \le 0, \quad j=1,2,\dots,J; \quad h_k(\mathbf{x}) = 0, \quad k=1,2,\dots,K
\end{equation}

\subsection{Dominasi Pareto dan Pareto Front}
\begin{tcolorbox}[colback=white,colframe=unibblue,title=\textbf{Konsep Dominasi Pareto}]
Suatu solusi Vektor keputusan $\mathbf{x}_a$ dikatakan **mendomina** solusi $\mathbf{x}_b$ ($\mathbf{x}_a \succ \mathbf{x}_b$) jika:
\begin{enumerate}
    \item $f_m(\mathbf{x}_a) \le f_m(\mathbf{x}_b)$ untuk semua $m \in \{1, 2, \dots, M\}$.
    \item Terdapat minimal satu fungsi tujuan $k$ di mana $f_k(\mathbf{x}_a) < f_k(\mathbf{x}_b)$.
\end{enumerate}
\end{tcolorbox}

Kumpulan seluruh solusi dalam $\Omega$ yang tidak mendomina satu sama lain membentuk **Pareto Optimal Set**, dan proyeksinya pada ruang fungsi tujuan disebut **Pareto Front**.

\section{Combined Economic dan Emission Dispatch (CEED)}

\subsection{Fungsi Tujuan 1: Biaya Bahan Bakar (\$/jam)}
\begin{equation}
    F_{\text{cost}}(\mathbf{P}_G) = \sum_{i=1}^{N_g} \left( a_i P_{G,i}^2 + b_i P_{G,i} + c_i \right)
\end{equation}

\subsection{Fungsi Tujuan 2: Emisi Gas Polutan (kg/jam)}
Emisi polutan $NO_x$ dari pembangkitan thermal dimodelkan sebagai:
\begin{equation}
    E_{\text{emission}}(\mathbf{P}_G) = \sum_{i=1}^{N_g} \left( \alpha_i P_{G,i}^2 + \beta_i P_{G,i} + \gamma_i \right)
\end{equation}

\subsection{Kendala Keseimbangan Daya dan Batas Pembangkit}
\begin{equation}
    \sum_{i=1}^{N_g} P_{G,i} - P_D - P_L = 0, \quad P_{i,\min} \le P_{G,i} \le P_{i,\max}
\end{equation}

\section{Adopsi Metode Deterministik dari Rekayasa Mesin}

\subsection{1. Quasi-Newton BFGS}
Metode BFGS (Broyden-Fletcher-Goldfarb-Shanno) memperbarui matriks taksiran Hessian $B_k \approx \nabla^2 f$ secara berulang tanpa membutuhkan komputasi turunan kedua:
\begin{equation}
    B_{k+1} = B_k + \frac{\mathbf{y}_k \mathbf{y}_k^T}{\mathbf{y}_k^T \mathbf{s}_k} - \frac{B_k \mathbf{s}_k \mathbf{s}_k^T B_k}{\mathbf{s}_k^T B_k \mathbf{s}_k}
\end{equation}
di mana $\mathbf{s}_k = \mathbf{x}_{k+1} - \mathbf{x}_k$ dan $\mathbf{y}_k = \nabla f(\mathbf{x}_{k+1}) - \nabla f(\mathbf{x}_k)$.

\subsection{2. Augmented Lagrangian Method}
Untuk menjamin kendala persamaan $h(\mathbf{x}) = 0$ terpenuhi secara presisi:
\begin{equation}
    L_A(\mathbf{x}, \lambda, \mu) = f(\mathbf{x}) + \lambda h(\mathbf{x}) + \frac{\mu}{2} [h(\mathbf{x})]^2
\end{equation}
Persamaan ini menggabungkan stabilitas penalti kuadratik dengan keakuratan pengali Lagrange $\lambda$.

\section{Penentuan Solusi Kompromi Terbaik (Fuzzy Logic)}
Untuk memilih 1 solusi operasional dari *Pareto Front*, digunakan fungsi keanggotaan linier *Fuzzy*:
\begin{equation}
    \mu_i^{(k)} = 
    \begin{cases}
        1, & f_i \le f_i^{\min} \\
        \frac{f_i^{\max} - f_i^{(k)}}{f_i^{\max} - f_i^{\min}}, & f_i^{\min} < f_i < f_i^{\max} \\
        0, & f_i \ge f_i^{\max}
    \end{cases}
\end{equation}
Solusi dengan nilai bobot keanggotaan ternormalisasi terbesar $\mu^{(k)}$ dipilih sebagai **Best Compromise Solution**.

\section{Implementasi Program Python}

\begin{lstlisting}[caption={Skrip Python Multi-Objective CEED dan Pareto Frontier}]
import numpy as np
import matplotlib.pyplot as plt

# Specs 3 Generator: [a, b, c, alpha, beta, gamma, Pmin, Pmax]
gen_specs = np.array([
    [0.00156, 7.92, 561, 0.00419, 0.3276, 13.86, 100, 600],
    [0.00194, 7.85, 310, 0.00419, 0.3276, 13.86, 100, 400],
    [0.00482, 7.97,  78, 0.00683, -0.545, 40.26,  50, 200]
])

P_D = 850.0  # MW

def calc_objectives(P1, P2):
    P3 = P_D - P1 - P2
    P = np.array([P1, P2, P3])
    
    # Cost ($/h)
    cost = np.sum(gen_specs[:, 0]*P**2 + gen_specs[:, 1]*P + gen_specs[:, 2])
    # Emission (kg/h)
    emission = np.sum(gen_specs[:, 3]*P**2 + gen_specs[:, 4]*P + gen_specs[:, 5])
    return cost, emission

print("Skrip simulasi Multi-Objective CEED siap dieksekusi.")
\end{lstlisting}

\section{Kesimpulan}
1. Pembelajaran hibrid yang mengawinkan kekuatan metaheuristik EE (PSO/GA) dan optimisasi deterministik Mesin (BFGS/Augmented Lagrangian) menghasilkan kemampuan analisis rekayasa tingkat tinggi.
2. Penggunaan *Pareto Frontier* dan *Fuzzy Decision Making* memberikan dasar ilmiah objektif bagi insinyur listrik dalam mengambil keputusan operasional *Smart Grid*.

\end{document}
